% !PRINTED 546
\weberpage{546}
%<<CONT
\[
\begin{aligned}
m &= 14,\ 34,\ 46,\ 82,\ 142 \equiv \pm2 \quad (\mathrm{mod}\ 16),\\
m &= 17,\ 49,\ 73,\ 97,\ 193 \equiv 1 \qquad (\mathrm{mod}\ 8).
\end{aligned}
\]
In these cases the classes of forms of the \gesperrt{principal genus} can be represented, for an even $m$, by

\weq{1}{(1, 0, m), \quad \left(2, 0, \dfrac{m}{2}\right),}

for an odd $m$ by

\weq{2}{(1, 0, m), \quad \left(2, 1, \dfrac{m+1}{2}\right),}

of which the latter is equivalent to

\weq{3}{\left(\dfrac{m+1}{2},\ \dfrac{m-1}{2},\ \dfrac{m+1}{2}\right).}

For the forms (1) the class invariants are, by § 127, 6.,
\[
\frac{f_1(\sqrt{-m})^2}{\sqrt2} \quad \text{and} \quad \frac{1}{\sqrt2} f_2\left(\frac{\sqrt{-m}}{2}\right)^2 = \frac{\sqrt2}{f_1(\sqrt{-m})^2} \qquad (\S\ 34)
\]
and for the forms (2), (3) (§ 127, 3.)
\[
\frac{f(\sqrt{-m})^2}{\sqrt2} \quad \text{and} \quad \frac{1}{\sqrt2} f\left(\frac{1+\sqrt{-m}}{1-\sqrt{-m}}\right)^2 = \frac{\sqrt2}{f(\sqrt{-m})^2},
\]
and by § 142 these are integral algebraic numbers.

If then we set, corresponding to the two cases:

\weq{4}{\sqrt2\, x = f_1(\sqrt{-m})^2, \quad f(\sqrt{-m})^2,}

then
\[
x + \frac{1}{x}
\]
is an integral algebraic number which contains only one square root, and from § 138 we obtain information as to which square root occurs in it.

It is $\sqrt2$ when $m \equiv 6$ $(\mathrm{mod}\ 8)$, hence for $m = 14, 46, 142$; it is $\sqrt{\dfrac{m}{2}}$ when $m \equiv 2$ $(\mathrm{mod}\ 8)$, hence for $m = 34, 82$; and further $\sqrt m$ in the case of an odd $m$ (with the exception of $m = 49$, where $\sqrt7$ takes its place).

If then we set

\weq{5}{x + \frac{1}{x} = a + b\sqrt p,}

then $a, b$ are rational numbers which have at most the denominator 2. But $a$ and $b$ must also be positive. For if in (5) we change the sign of $\sqrt p$, there arises another quadratic
%CONT>>
